\makeabstract
{
Is Nye Theory Accurate or Antiquated: Verifying Modeled Crevasse Depths Against Radar-Derived Crevasse Depths in Pine Island and Thwaites Glaciers
}
{
Emily Lan (1), Nick Holschuh, PhD (1)
}
{
\textsuperscript{1}Department of Geology, Amherst College, Amherst, MA, USA
}
{
There has been no comprehensive testing of the accuracy of Nye's theory of crevassing in predicting crevasse formation and depth. This study models basal crevasse depths according to the principles of Nye and corroborates them against manually derived crevasse depths from radar data in the high strain-rate Pine Island and Thwaites glaciers.
}
{
Using the Nye crevasse depth formula and using skin temperature, ice thickness, and velocity datasets for our calculations yielded our map for expected crevasse depths. Filtering through the available radar data for where the flight path is approximately parallel to ice flow direction, we collected crevasse tip and ice bottom depths, from which we determined the observed basal crevasse depths.
}
{
We found that the true values of basal crevasse depth range from 1.64x to 2.52x the values predicted by our model. We found that less than 20\% of the variability in observed values was explained by Nye’s model of crevassing when using either the skin temperature or melting point as temperature and the model only accurately predicted 6\% of basal crevasse depths.
}
{
This study provides strong evidence that Nye theory underestimates the depth of basal crevasses. A valuable question for further study is whether this underestimation could have been the result of incorrect values for the creep parameter A. In conclusion, it may be more appropriate to use newer models such as Linear Elastic Fracture Mechanics (LEFM) for use in predicting glacier system behavior.
}

\newpage
